\chapter{INTRODUCTION}
\section{Overview}
\pagenumbering{arabic}

A sequence model performs \emph{state tracking} when it maintains a structured
internal quantity whose correct value at position $t$ depends on every symbol
observed before it. Tracking the configuration of a board across a sequence of
moves, the value of a variable across a program trace, or the orientation of a
rigid body across a stream of rotations are all instances of the same
computational pattern. Attention-based architectures address it by re-reading the
entire history at every position, which costs time quadratic in sequence length
\cite{attention}. Linear recurrent architectures, including linear attention
\cite{lineartransformers}, structured state-space models \cite{s4, mamba}, and
the delta-rule family \cite{fastweight, deltanet}, instead carry a fixed-size
state forward and therefore achieve cost linear in length. This efficiency has
historically come at the expense of exactly the capability described above:
models whose state-transition operator is diagonal are provably unable to track
non-commutative structure \cite{illusionofstate, diagonalssm}.

Permutation-group word problems are the standard probe for this capability. The
model receives a sequence of group elements and must emit, at every position, the
cumulative product of all elements seen so far. The task admits no shortcut, its
correct answer is computable exactly, and its difficulty can be varied in a
controlled manner by changing the alphabet of group elements from which sequences
are drawn.

The DeltaProduct layer \cite{deltaproduct} narrows the gap between linear cost
and state-tracking capability by constructing each state-transition matrix as a
product of $n_h$ generalized Householder factors,
\begin{equation}
A_t \;=\; \prod_{i=1}^{n_h} \bigl( I - \beta_{t,i}\, k_{t,i} k_{t,i}^{\!\top} \bigr),
\qquad \lVert k_{t,i} \rVert = 1, \quad \beta_{t,i} \in [0,2].
\label{eq:deltaproduct}
\end{equation}
Each factor is a rank-one modification of the identity, and the number of such
factors controls the expressivity of the set of transitions the layer can
realise. Increasing $n_h$ enlarges that set and increases computation in direct
proportion. In the formulation as published, $n_h$ is a single global
hyperparameter: it is chosen once, before training, and the same number of
factors is applied to every token of every sequence.

This work is concerned with that design decision. It establishes what the correct
per-token factor count is, shows that this quantity cannot be read off individual
tokens, and develops a layer in which the count is predicted per token by a
learned halting gate trained under a penalty on expected computation.

\section{Motivation}

The immediate argument for a per-token factor count is one of efficiency. On a
sequence in which most symbols require a single factor and a minority require
four, a fixed-order model must be provisioned for the maximum and therefore pays
that maximum everywhere. Allocating per token reduces the required work to the
expected requirement rather than its supremum, and the reduction is larger the
more skewed the distribution of requirements.

A second and stronger argument follows from the analysis developed in this work.
The number of Householder factors that a symbol requires is not a property of the
symbol. It is a property of the algebraic closure of the entire alphabet from
which symbols are drawn. A layer is free to realise any faithful representation
of the group generated by the alphabet, and the factor count required for a given
element is the rank of that element's deviation from the identity, minimised over
all such representations. The consequence is concrete and counter-intuitive:
adding a single transposition to an alphabet of five-cycles changes the generated
group from the alternating group $A_5$ to the symmetric group $S_5$, removes
access to the three-dimensional faithful representation that $A_5$ admits, and
thereby doubles the requirement of every five-cycle in the alphabet, from two
factors to four. The least expressive symbol available makes every other symbol
more expensive.

This has a practical implication for the hyperparameter itself. Selecting $n_h$
correctly requires knowing the minimum faithful representation dimension of the
group generated by the data, a quantity that is not available from inspection of
a dataset and that changes discontinuously when the alphabet changes. A
practitioner who fixes $n_h$ globally must therefore over-provision, wasting
computation on every token, or under-provision, in which case the layer cannot
represent the required transitions and the failure is silent. A layer that infers
its own order removes a hyperparameter that cannot be set reliably by hand, and
this argument is independent of any efficiency gain.

A third consideration concerns measurement. The training behaviour of these
architectures on the harder alphabets is bimodal rather than concentrated: a run
either escapes an extended loss plateau and solves the task, or remains at chance
accuracy, with intermediate outcomes rare. Conclusions drawn from single training
runs on such configurations are therefore unreliable, and the experimental
protocol adopted here reports escape fractions over repeated runs in place of
point estimates wherever this behaviour is present.

\section{Contributions}

The contributions of this work are the following.

\begin{enumerate}[label={(\arabic*)}]
	\item An order-requirement model for Householder-product state transitions,
	establishing that the factor count required by a symbol is determined by the
	minimum faithful representation dimension of the group generated by the whole
	alphabet, together with its verification by explicit construction in double
	precision.
	\item Two reachability bounds for products of generalized Householder factors:
	a rank bound relating the factor count to the transposition length of the
	target permutation, and a determinant-parity bound that determines when a
	shorter solution may be padded to a longer factor count. The second bound
	constrains the design of any gating mechanism applied to such a layer.
	\item An adaptive-order layer in which a monotone halting gate predicts a
	per-token factor count, trained end to end with a penalty on expected
	computation and without supervision on the allocation.
	\item An empirical characterisation of capability at matched computation,
	quantifying the reduction in required work that a correct per-token allocation
	achieves relative to the cheapest adequate fixed order.
	\item A stability analysis over fifty training runs establishing the bimodal
	character of training outcomes on these tasks, and the measurement protocol
	that follows from it.
\end{enumerate}
